\cvsection{Expériences professionnelles}

\begin{cventries}
  \cventry
    {Stagiaire ingénieur de recherche en IA}
    {INRIA}
    {Bordeaux, France}
    {Mars 2026 -- Septembre 2026}
    {
      \begin{cvitems}
        \item {Conception et développement en autonomie d’une application web de bout en bout intégrant un agent IA pour la recherche d’articles scientifiques, à partir des besoins définis avec mon tuteur.}
        \item {Développement de l’agent en Python sans framework d’orchestration, avec intégration de l’API OpenAI, appels d’outils et gestion de l’état, de la mémoire et du contexte.}
        \item {Conception de prompts pour adapter le comportement de l’agent au cas d’usage scientifique et guider l’utilisation de ses outils.}
        \item {Mise en œuvre d’un pipeline RAG avec chunking, embeddings et recherche sémantique dans ChromaDB, associé à une persistance MongoDB.}
        \item {Développement du backend de l’application en Python/FastAPI, avec API REST documentées avec OpenAPI et stockage des documents et métadonnées dans MongoDB.}
        \item {Développement de l’interface React/Next.js pour interagir avec l’agent et ses outils, avec sélection par l’utilisateur des messages à inclure dans le contexte du modèle.}
        \item {Automatisation en Python de la collecte de métadonnées d’articles scientifiques en ligne et transmission des résultats de tâches cron par webhooks.}
        \item {Évaluation de plusieurs LLM et quantifications selon la qualité, la latence et le contexte, analyse des résultats et rédaction en cours d’un article scientifique.}
        \item {Conteneurisation de composants avec Docker, versionnement sous Git et mise en place d’une CI/CD sur GitHub.}
        \item {Évaluation d’OpenClaw dans un environnement Docker isolé et analyse des risques liés à l’accès autonome aux outils et aux fichiers.}
      \end{cvitems}
    }

  \cventry
    {Stagiaire ingénieur de recherche en apprentissage par renforcement}
    {Université Karunya}
    {Coimbatore, Inde}
    {Juin 2025 -- Sept. 2025}
    {
      \begin{cvitems}
        \item {Développement d’un modèle d’apprentissage par renforcement et de son environnement d’entraînement pour la navigation autonome d’un robot hospitalier.}
        \item {Déploiement et optimisation du logiciel et du modèle sur Raspberry Pi pour des tests en conditions réelles.}
      \end{cvitems}
    }
\end{cventries}
